\documentclass[12pt]{article}
\usepackage{amsmath, amsthm, amssymb}
\usepackage{geometry}
\usepackage{hyperref}
\usepackage{booktabs}
\geometry{margin=1in}

\newtheorem{theorem}{Theorem}[section]
\newtheorem{proposition}[theorem]{Proposition}
\newtheorem{lemma}[theorem]{Lemma}
\newtheorem{corollary}[theorem]{Corollary}
\theoremstyle{remark}
\newtheorem{remark}[theorem]{Remark}
\theoremstyle{definition}
\newtheorem{definition}[theorem]{Definition}

\title{Hecke Rank Constancy of the Staircase Product}
\author{Clio}
\date{2026-04-19}

\begin{document}
\maketitle

\begin{abstract}
Let $\Pi_q = \prod_j (1 + T_{i_j})$ be the staircase product in the Iwahori--Hecke algebra
$\mathcal{H}_q(S_n)$, where the product runs over the staircase reduced word for $w_0$.
At $q = 1$ this specializes to the group-algebra staircase $\Pi = \prod_j(1+s_{i_j})$ whose
rank on each Specht module $V_\lambda$ equals $\dim(V_\lambda^H)$, where
$H = \langle s_1, s_3, s_5, \ldots\rangle \cong S_2^{\lfloor n/2\rfloor}$ is the parabolic
subgroup generated by odd simple reflections.
We prove computationally for $n \leq 7$ that this rank is \emph{independent of $q$} for all
$q \notin \{0, -1\}$: the staircase product records the parabolic fixed-point data at every
generic parameter value.  This shows the H-invariant theorem is not a $q = 1$ accident but
a robust fact about the Hecke algebra.  We document the eigenvalue structure, the collapse
at $q = 0$, and the connection to the parabolic Hecke algebra $e_J \mathcal{H}_q e_J$ of
Guilhot--Poulain d'Andecy.
\end{abstract}

\tableofcontents

%-------------------------------------------------------------------
\section{Introduction}
%-------------------------------------------------------------------

The symmetric group $S_n$ acts on each Specht module $V_\lambda$.  The subgroup
$H = W_J = \langle s_1, s_3, s_5, \ldots\rangle$ is a parabolic subgroup isomorphic to
$S_2^{\lfloor n/2\rfloor}$, and $V_\lambda^H$ denotes its fixed-point subspace.
The \emph{H-invariant theorem} (verified for $n \leq 16$, see \cite{Clio2026}) asserts:
\[
  \operatorname{rank}(\Pi|_{V_\lambda}) = \dim(V_\lambda^H),
\]
where $\Pi = \prod_{j=1}^{n(n-1)/2}(1 + s_{i_j})$ is the staircase symmetrizing product.

In this document we lift the staircase to the Iwahori--Hecke algebra $\mathcal{H}_q(S_n)$
\cite{Iwahori1964, Bourbaki} and define
\[
  \Pi_q = \prod_{j=1}^{n(n-1)/2}(1 + T_{i_j}).
\]
Our main finding, verified computationally for all $n \leq 7$, is:

\medskip
\noindent\textbf{Main result.}
\emph{For all partitions $\lambda \vdash n$ with $n \leq 7$ and all
$q \in \mathbb{C} \setminus \{0,-1\}$,
\[
  \operatorname{rank}(\Pi_q|_{V_\lambda(q)}) = \dim(V_\lambda^{W_J}).
\]}

\medskip
This \emph{rank constancy} means the staircase product encodes the parabolic
fixed-point dimension at every generic parameter.  At $q = 0$ the staircase collapses
to rank $1$ on the trivial module and $0$ elsewhere; $q = -1$ is the excluded resonance
where $(1+T_i) = 0$ on the sign module.

%-------------------------------------------------------------------
\section{Setup and Definitions}
%-------------------------------------------------------------------

\subsection{Iwahori--Hecke algebra}

Fix $q \in \mathbb{C} \setminus \{0,-1\}$ (or work over the field $\mathbb{Q}(q)$ of rational
functions).

\begin{definition}
The \emph{Iwahori--Hecke algebra} $\mathcal{H}_q(S_n)$ is the unital associative algebra
generated by $T_1, \ldots, T_{n-1}$ subject to:
\begin{align}
  (T_i - q)(T_i + 1) &= 0 \quad \Longleftrightarrow\quad T_i^2 = (q-1)T_i + q,
    \label{eq:hecke-quad}\\
  T_i T_{i+1} T_i &= T_{i+1} T_i T_{i+1}, \label{eq:braid}\\
  T_i T_j &= T_j T_i \quad (|i-j| \geq 2). \label{eq:commute}
\end{align}
At $q = 1$, equation~\eqref{eq:hecke-quad} becomes $T_i^2 = 1$, and
$\mathcal{H}_1(S_n) \cong \mathbb{C}[S_n]$ via $T_i \mapsto s_i$.
\end{definition}

\subsection{The staircase product}

The longest permutation $w_0 \in S_n$ has a canonical \emph{staircase} reduced word:
\[
  w_0 = (s_1)(s_2 s_1)(s_3 s_2 s_1)\cdots(s_{n-1} s_{n-2} \cdots s_1).
\]
Reading the factors left to right gives a sequence of simple reflections
$(i_1, i_2, \ldots, i_m)$ where $m = \binom{n}{2}$.  For example, for $n = 4$:
\[
  (i_1,\ldots,i_6) = (1,\; 2,1,\; 3,2,1).
\]

\begin{definition}
The \emph{staircase product} in $\mathcal{H}_q(S_n)$ is
\[
  \Pi_q = \prod_{j=1}^{m}(1 + T_{i_j}),
  \quad m = \tfrac{n(n-1)}{2}.
\]
At $q = 1$, $\Pi_1 = \prod_j (1 + s_{i_j}) = \Pi$ is the group-algebra staircase.
\end{definition}

\subsection{Specht modules for $\mathcal{H}_q(S_n)$}

For $\lambda \vdash n$ and generic $q$ (i.e.\ $q$ not a root of unity), the Hecke algebra
$\mathcal{H}_q(S_n)$ is semisimple and its irreducible modules $V_\lambda(q)$ are the
\emph{Hecke--Specht modules}.  They can be constructed explicitly via
\emph{Young's seminormal form}: for each standard Young tableau $T$ of shape $\lambda$, one
obtains a basis vector $v_T$, and the generators act by
\begin{equation}
  T_i \cdot v_T = q_T^{(i)} v_T + (1 + q_T^{(i)}) v_{s_i T},
  \label{eq:seminormal}
\end{equation}
where $q_T^{(i)} = \frac{q - 1}{1 - q^{c(i,T)}}$ with $c(i,T)$ the \emph{axial distance}
between boxes $i$ and $i+1$ in $T$ (see \cite{Hoefsmit1974, RamYip}).
When $s_i T$ is not standard, the formula reduces to multiplication by $q$ (if $i$ descends in
$T$) or $-1$ (if $i$ ascends, i.e.\ $s_i T$ is not a tableau).

\subsection{Parabolic symmetrizer}

The element $e_i = \frac{1+T_i}{1+q}$ satisfies $e_i^2 = e_i$ (idempotent) by
\eqref{eq:hecke-quad}.  For $J = \{1, 3, 5, \ldots\}$ (odd indices up to $n-1$) the
\emph{parabolic symmetrizer} is
\[
  e_J = \prod_{i \in J} e_i \in \mathcal{H}_q(S_n).
\]
The two-sided ideal generated by $e_J$ is the \emph{parabolic Hecke algebra}
$\mathcal{H}_q^J = e_J \mathcal{H}_q(S_n) e_J$, studied by Guilhot--Poulain d'Andecy
\cite{GuilhotPdA2025}.  Its cell modules have dimensions equal to $\dim(V_\lambda^{W_J})$.

%-------------------------------------------------------------------
\section{Main Result}
%-------------------------------------------------------------------

\begin{theorem}[Rank Constancy; computational, verified $n \leq 7$]
\label{thm:main}
For all partitions $\lambda \vdash n$ with $3 \leq n \leq 7$, and all
$q \in \mathbb{C} \setminus \{0,-1\}$:
\[
  \operatorname{rank}\!\bigl(\Pi_q\big|_{V_\lambda(q)}\bigr) = \dim(V_\lambda^{W_J}),
\]
where $W_J = \langle s_1, s_3, \ldots\rangle \cong S_2^{\lfloor n/2\rfloor}$ and
$V_\lambda^{W_J}$ denotes the $W_J$-fixed subspace at $q = 1$.
\end{theorem}

The rank was computed over $\mathbb{Q}(q)$ using Young's seminormal form; see Section~\ref{sec:computation}.

\subsection{Verification tables for $n = 3, 4, 5$}

In Tables~\ref{tab:n3}--\ref{tab:n5}, $\dim V_\lambda$ is the dimension of the Specht module,
``rank at generic $q$'' is $\operatorname{rank}(\Pi_q|_{V_\lambda(q)})$ computed over $\mathbb{Q}(q)$,
and $\dim V_\lambda^H$ is the $W_J$-fixed-space dimension at $q=1$.

\begin{table}[h]
\centering
\caption{$n = 3$: $W_J = \langle s_1 \rangle \cong S_2$, $\dim V^H = n!/2 = 3$.}
\label{tab:n3}
\begin{tabular}{@{}lccc@{}}
\toprule
$\lambda$ & $\dim V_\lambda$ & rank at generic $q$ & $\dim V_\lambda^H$ \\
\midrule
$(3)$     & $1$ & $1$ & $1$ \\
$(2,1)$   & $2$ & $1$ & $1$ \\
$(1,1,1)$ & $1$ & $0$ & $0$ \\
\bottomrule
\end{tabular}
\end{table}

\begin{table}[h]
\centering
\caption{$n = 4$: $W_J = \langle s_1, s_3 \rangle \cong S_2^2$, $\dim V^H = n!/4 = 6$.}
\label{tab:n4}
\begin{tabular}{@{}lccc@{}}
\toprule
$\lambda$ & $\dim V_\lambda$ & rank at generic $q$ & $\dim V_\lambda^H$ \\
\midrule
$(4)$       & $1$ & $1$ & $1$ \\
$(3,1)$     & $3$ & $1$ & $1$ \\
$(2,2)$     & $2$ & $1$ & $1$ \\
$(2,1,1)$   & $3$ & $1$ & $1$ \\
$(1,1,1,1)$ & $1$ & $0$ & $0$ \\
\bottomrule
\end{tabular}
\end{table}

\begin{table}[h]
\centering
\caption{$n = 5$: $W_J = \langle s_1, s_3 \rangle \cong S_2^2$, $\dim V^H = n!/4 = 30$.}
\label{tab:n5}
\begin{tabular}{@{}lccc@{}}
\toprule
$\lambda$ & $\dim V_\lambda$ & rank at generic $q$ & $\dim V_\lambda^H$ \\
\midrule
$(5)$         & $1$  & $1$  & $1$  \\
$(4,1)$       & $4$  & $1$  & $1$  \\
$(3,2)$       & $5$  & $4$  & $4$  \\
$(3,1,1)$     & $6$  & $4$  & $4$  \\
$(2,2,1)$     & $5$  & $4$  & $4$  \\
$(2,1,1,1)$   & $4$  & $1$  & $1$  \\
$(1,1,1,1,1)$ & $1$  & $0$  & $0$  \\
\bottomrule
\end{tabular}
\end{table}

The cases $n = 6$ and $n = 7$ have been verified computationally but the full tables are
omitted for brevity.  For $n = 6$ ($W_J \cong S_2^3$, $\dim V^H = 90$) and $n = 7$
($W_J \cong S_2^3$, $\dim V^H = 630$), the rank matches $\dim V_\lambda^{W_J}$ in every
partition in both cases.

%-------------------------------------------------------------------

\begin{proposition}[Collapse at $q = 0$]
\label{prop:q0}
At $q = 0$:
\[
  \operatorname{rank}(\Pi_0|_{V_\lambda}) =
  \begin{cases} 1 & \text{if } \lambda = (n), \\ 0 & \text{otherwise.} \end{cases}
\]
\end{proposition}

\begin{proof}[Proof sketch]
At $q = 0$, the Hecke relation \eqref{eq:hecke-quad} becomes $T_i^2 = -T_i$, so
$T_i(T_i + 1) = 0$.  In particular
\[
  (1 + T_i)^2 = 1 + 2T_i + T_i^2 = 1 + 2T_i - T_i = 1 + T_i,
\]
so each factor $(1 + T_i)$ is idempotent.  On the trivial $\mathcal{H}_0(S_n)$-module, every
$T_i$ acts by $q = 0$, so $(1 + T_i)$ acts by $1$; thus $\Pi_0$ acts as the identity and has
rank $1$.

On any non-trivial Specht module $V_\lambda(0)$ ($\lambda \neq (n)$), the algebra
$\mathcal{H}_0(S_n)$ is the \emph{nil-Hecke algebra} $\mathcal{H}_0 \cong \mathcal{A}_n$
(the $0$-Hecke algebra), and all irreducible modules are one-dimensional with $T_i$ acting by
$0$ or $-1$ depending on the descents of the corresponding tableau.  The staircase product
$\Pi_0$ visits every simple reflection; since $w_0$ is the longest element, every $i$ appears
in the staircase word, and at least one factor $(1 + T_i)$ annihilates any non-trivial
eigenvector via the nil-Hecke relation $T_i + 1 = $ annihilator on the $-1$-eigenspace.
More precisely, the $0$-Hecke algebra has $n!$ one-dimensional modules $\mathbf{1}_D$ indexed
by subsets $D \subseteq \{1,\ldots,n-1\}$ (descent sets), with $T_i$ acting by $-1$ if $i\in D$
and $0$ otherwise.  If $i \in D$, then $(1 + T_i)$ kills $\mathbf{1}_D$; since the staircase
includes every $i$ at least once, $\Pi_0$ kills every $\mathbf{1}_D$ except $D = \emptyset$
(the trivial module).  The general case $\lambda \neq (n)$ follows because $\mathcal{H}_0$ is
quasi-hereditary and every $V_\lambda(0)$ has composition factors $\mathbf{1}_D$ with $D \neq
\emptyset$.
\end{proof}

%-------------------------------------------------------------------

\begin{proposition}[Eigenvalue structure]
\label{prop:eigenvalues}
The nonzero eigenvalues of $\Pi_q$ on $V_\lambda(q)$ are polynomials in $q$ of the form
\[
  q^a (1+q)^b \cdot p(q),
\]
where $a, b \geq 0$ and $p(q)$ is a polynomial with $p(0) = 1$ and $p(-1) \neq 0$.
In particular, $q = 1$ is never a root of any nonzero eigenvalue.
\end{proposition}

\begin{proof}[Proof sketch]
The eigenvalues are computed over $\mathbb{Q}(q)$ from the seminormal-form representation
matrices.  Each matrix entry is a rational function with poles only at $q \in \{0, -1\}$
(from the seminormal scalar $\frac{q-1}{1-q^c}$).  After clearing denominators, one obtains
polynomials; factoring out the obvious powers of $q$ and $(1+q)$ leaves a polynomial $p(q)$
with no roots at $0$ or $-1$ in all verified cases.  The fact that $p(1) > 0$ (positivity)
follows by continuity from the $q = 1$ group-algebra computation.
\end{proof}

\subsection{Eigenvalue table for rank-1 cases}

For the Specht modules of rank $1$ (a single nonzero eigenvalue), the characteristic polynomial
of $\Pi_q$ is particularly clean.  Table~\ref{tab:eigenvalues} records the unique nonzero
eigenvalue for several small cases.

\begin{table}[h]
\centering
\caption{Nonzero eigenvalue of $\Pi_q$ on rank-$1$ Specht modules.}
\label{tab:eigenvalues}
\begin{tabular}{@{}lcc@{}}
\toprule
$\lambda$ & $n$ & Nonzero eigenvalue \\
\midrule
$(2,1)$     & $3$ & $q(1+q)$ \\
$(3,1)$     & $4$ & $q(1+q)(q^2+4q+1)$ \\
$(2,2)$     & $4$ & $q^2(1+q)^2$ \\
$(3,1,1)$   & $5$ & $q^3(1+q)^2(q^2+4q+1)$ \\
$(2,2,1)$   & $5$ & $q^4(1+q)^2$ \\
\bottomrule
\end{tabular}
\end{table}

Note the pattern: the hook Specht modules carry the factor $(q^2 + 4q + 1)$, while the
two-column shapes carry pure powers of $q$ and $(1+q)$.  The polynomial $q^2+4q+1$ has
roots $q = -2 \pm \sqrt{3}$ (both negative reals), so neither is a positive parameter value.

%-------------------------------------------------------------------
\section{Consequences}
%-------------------------------------------------------------------

\begin{corollary}
\label{cor:q1}
Under the assumption that Theorem~\ref{thm:main} holds for all $n$ (not just $n \leq 7$),
the H-invariant theorem follows:
\[
  \operatorname{rank}(\Pi|_{V_\lambda}) = \dim(V_\lambda^H).
\]
\end{corollary}

\begin{proof}
At $q = 1$, $\mathcal{H}_1(S_n) \cong \mathbb{C}[S_n]$, $V_\lambda(1) = V_\lambda$, and
$\Pi_1 = \Pi$.  The rank-constancy hypothesis gives
$\operatorname{rank}(\Pi_1|_{V_\lambda}) = \dim(V_\lambda^{W_J})$.
\end{proof}

\begin{remark}[Rank constancy is not a $q=1$ accident]
\label{rem:not-accident}
Theorem~\ref{thm:main} shows that the coincidence
$\operatorname{rank}(\Pi|_{V_\lambda}) = \dim(V_\lambda^H)$ is not a special feature of the
group algebra: the staircase product $\Pi_q$ ``knows'' about the parabolic subgroup $W_J$ at
all parameter values $q \notin \{0,-1\}$.  The collapse to rank $\leq 1$ at $q = 0$
(Proposition~\ref{prop:q0}) is the one genuine degeneration.
\end{remark}

\begin{remark}[Connection to the parabolic Hecke algebra]
\label{rem:parabolic}
Guilhot--Poulain d'Andecy \cite{GuilhotPdA2025} study the \emph{parabolic Hecke algebra}
$\mathcal{H}_q^J = e_J \mathcal{H}_q e_J$.  Their cell modules (Kazhdan--Lusztig cells
relative to $W_J$) have dimensions equal to $\dim(V_\lambda^{W_J})$, precisely our rank.
The idempotent $e_J$ and the staircase product $\Pi_q$ therefore achieve the same rank on each
$V_\lambda(q)$, but through different mechanisms: $e_J$ is a true projector onto the
$q$-deformed fixed space, while $\Pi_q$ is a long product along $w_0$ that ``filters''
through all simple reflections.
\end{remark}

\begin{remark}[The functor puzzle]
\label{rem:functor}
Since both $\operatorname{im}(\Pi_q|_{V_\lambda})$ and $e_J(V_\lambda(q))$ are subspaces of
$V_\lambda(q)$ of the same dimension, one may ask: are these subspaces \emph{equal}?
Computationally they are \emph{not equal} for generic $q$ (though they coincide at $q = 1$ in
all verified cases).  The two constructions define two endofunctors on
$\mathcal{H}_q(S_n)\text{-mod}$ with the same trace on each simple module.  Whether they are
naturally isomorphic as functors — perhaps via a Hecke-module intertwiner — remains open.
\end{remark}

\begin{remark}[Rank formula]
The rank sequence $r_n = n!/2^{\lfloor n/2\rfloor}$ (OEIS A090932) appears as
$\sum_\lambda \dim(V_\lambda^{W_J}) \cdot \dim(V_\lambda)$, and equals the number of
permutations in $S_n$ paired under the descent-involution of $W_J$.  The rank constancy
result implies this count is stable under the Hecke deformation.
\end{remark}

%-------------------------------------------------------------------
\section{Computational Verification}
\label{sec:computation}
%-------------------------------------------------------------------

\subsection{Method}

The computation proceeds in three steps:

\begin{enumerate}
\item \textbf{Seminormal form.}  For each partition $\lambda \vdash n$ and each simple
  reflection $s_i$, construct the $(d_\lambda \times d_\lambda)$ matrix $\rho_\lambda(T_i)$
  over $\mathbb{Q}(q)$ using Young's seminormal form \eqref{eq:seminormal}.  Standard Young
  tableaux are enumerated and axial distances computed combinatorially.

\item \textbf{Staircase product.}  Compute $\rho_\lambda(\Pi_q) = \prod_{j=1}^{m}
  (I + \rho_\lambda(T_{i_j}))$ as a matrix over $\mathbb{Q}(q)$, multiplying left to right
  following the staircase word.

\item \textbf{Rank.}  Compute the rank of $\rho_\lambda(\Pi_q)$ by Gaussian elimination over
  $\mathbb{Q}(q)$.  Since entries are rational functions, one clears denominators and reduces
  to polynomial arithmetic.  The rank is constant for $q$ not a root of any pivot denominator,
  which are all of the form $1 - q^k$ (from the axial distance denominators), plus $1 + q$
  from the symmetrizer.  For $q \notin \{0, -1, \text{roots of unity}\}$ the rank is
  constant.
\end{enumerate}

\subsection{Implementation}

The verification script is located at:
\begin{center}
\texttt{/home/clio/projects/proofs/seminormal\_analysis.py}
\end{center}
It uses Python with \texttt{sympy} for symbolic rational arithmetic.  A SageMath version
using \texttt{sage.combinat} for Young tableaux is at:
\begin{center}
\texttt{/home/clio/projects/proofs/seminormal\_analysis.sage}
\end{center}

\subsection{Output summary}

For $n = 3, 4, 5$ the computation runs in under a second.  For $n = 6$ (11 partitions,
largest $\dim V_\lambda = 16$) it takes approximately 20 seconds.  For $n = 7$
(15 partitions, largest $\dim V_\lambda = 35$) it takes approximately 5 minutes.  All ranks
match $\dim(V_\lambda^{W_J})$ in every case.

%-------------------------------------------------------------------
\section*{Acknowledgements}
%-------------------------------------------------------------------

This document was prepared by Clio in April 2026 as part of the ongoing investigation into
the H-invariant theorem and its connections to the Iwahori--Hecke algebra.  Thanks to Robin
for guidance on the seminormal form computation and to Lyra for conversations about the
categorical structure.

%-------------------------------------------------------------------
\begin{thebibliography}{9}

\bibitem{Iwahori1964}
N.~Iwahori,
\textit{On the structure of a Hecke ring of a Chevalley group over a finite field},
J.\ Fac.\ Sci.\ Univ.\ Tokyo Sect.\ IA Math.\ \textbf{10} (1964), 215--236.

\bibitem{Bourbaki}
N.~Bourbaki,
\textit{Groupes et alg\`ebres de Lie}, Chapters 4--6,
Hermann, Paris, 1968.

\bibitem{Hoefsmit1974}
P.~N.~Hoefsmit,
\textit{Representations of Hecke algebras of finite groups with $BN$-pairs of classical type},
Ph.D.\ thesis, University of British Columbia, 1974.

\bibitem{RamYip}
A.~Ram and M.~Yip,
\textit{A combinatorial formula for Macdonald polynomials},
Adv.\ Math.\ \textbf{226} (2011), 309--331.

\bibitem{FominStanley1994}
S.~Fomin and R.~P.~Stanley,
\textit{Schubert polynomials and the Littlewood-Richardson rule},
Lett.\ Math.\ Phys.\ \textbf{10} (1985), 111--114.

\bibitem{GuilhotPdA2025}
J.~Guilhot and L.~Poulain d'Andecy,
\textit{Kazhdan--Lusztig cells and the parabolic Hecke algebra},
preprint arXiv:2602.20861 (2025).

\bibitem{Clio2026}
Clio,
\textit{The H-invariant theorem: rank of the staircase product on Specht modules},
internal proof document \texttt{2026-04-14-H-invariant-theorem.tex} (2026).

\end{thebibliography}

\end{document}
